\documentclass[11pt]{letter}

\usepackage[margin=0.5in]{geometry}
\usepackage{hyperref}

\begin{document}

\begin{letter}{}

\opening{Dear Hiring Team,}

I am writing to apply for the New Grad Software Engineer position at Magical in Toronto. With a Master of Science in Computer Science from the University of Calgary and hands-on experience building backend APIs, Python/C++ data systems, and LLM-enabled automation workflows, I am excited by the chance to help build reliable AI agents for complex healthcare operations.

What stands out to me about Magical is that the role is not just about wrapping a model in a product interface. The work involves long-running workflows, state, queues, browser automation, model calls, observability, human review, and recovery from messy real-world edge cases. That is exactly the kind of engineering problem I enjoy: making an intelligent system useful, debuggable, and dependable when the input is imperfect and the path forward is not always obvious.

My backend experience started as a back-end developer at Asr Gooyesh Pardaz, where I built Django REST APIs for NLP, text-to-speech, and speech-to-text services backed by PostgreSQL. I implemented authentication flows, payment-service logic, and plan-management features, and contributed to a WebRTC-based chatbot system using GPT-style responses and FAQ data. This gave me a practical foundation in APIs, databases, user-facing backend logic, and debugging across services.

More recently, I have been building an agentic job-search automation platform using LLM workflows, n8n, OpenClaw, Docker, Google Sheets, Google Drive, SearXNG, and Telegram API integrations. The project involves prompt-driven analysis, structured outputs, workflow orchestration, spreadsheet tracking, and automation around real decisions. While it is not a healthcare product, it is closely aligned with Magical's need for engineers who can connect AI workflows, reason about model outputs, improve reliability, and build tools that support practical human decisions.

My graduate research also strengthened the engineering habits needed for high-ownership product work. With BigGeo, Mitacs, and the University of Calgary, I built Python and C++ pipelines for large-scale data ingestion, preprocessing, storage, retrieval, and evaluation. I improved one C++ processing workflow from 233 seconds to 26 seconds for a 10-tensor workload and reduced storage by 38\% on a scaled Canada-wide dataset. That work taught me to care about correctness, performance, documentation, and the small implementation details that make systems trustworthy.

I have not worked directly in healthcare automation before, and my professional TypeScript/React experience is lighter than my Python and backend experience. However, I bring strong software fundamentals, real backend development experience, hands-on AI automation work, research-level persistence, and a genuine interest in learning quickly in an early team. I would be excited to contribute to Magical's engineering work and grow into the kind of product-minded engineer who can own meaningful systems from implementation through monitoring and iteration.

Thank you for considering my application. I would welcome the opportunity to discuss how my backend, AI automation, data systems, and research software experience can contribute to Magical's AI agent platform.

\closing{Sincerely,

Amir Mirzai Golpayegani

\href{mailto:amirgolpaa24@gmail.com}{amirgolpaa24@gmail.com}}

\end{letter}

\end{document}
